\section{The remaining automorphic hypotheses on the marked host}

% BEGIN O72H1
\paragraph{O72H1. The coefficient representation and conjugation of a cusp.}
Let $E$ be any totally complex number field, let $n>0$, and suppose
the cusp $\pi$ has infinitesimal characters $m_\tau$ opposite to
those of the given irreducible algebraic representation $\Xi$. The
Harish--Chandra parameter of a highest weight $\lambda$ is the
Weyl orbit of $\lambda+\rho$. In particular its entries are
distinct and lie in $(n-1)/2+\mathbb Z$. The negative multiset
is the parameter of the dual. Thus $\pi$ is regular algebraic.
By Clozel, \emph{Motifs et formes automorphes}, \S3, a regular
algebraic cuspidal representation of $\mathrm{GL}_n$ has nonzero
relative Lie algebra cohomology for an irreducible algebraic
coefficient $\Xi'$. Wigner's lemma forces its dual infinitesimal
character to be $m_\tau$ for every embedding: a central element
acts on an Ext group alike from the two modules, so distinct
characters give zero Ext. An infinitesimal character determines an
irreducible algebraic representation at every embedding by taking
the dominant representative and subtracting $\rho$. Consequently
$\Xi'\simeq\Xi$. In the compact model the cochains are alternating
forms on the Hermitian tangent directions with values in
$\pi\otimes\Xi$, equivariant for the product of unitary groups.
The bracket of two tangent directions lies in the compact Lie algebra;
the differential is the alternating sum of their actions. This
identifies the nonzero class with the source's $\mathrm{IsCohFor}$.
A nonzero finite tensor factor may be fixed throughout. This gives
$\mathrm{cohC}$ for the specified $\Xi$.

For $\mathrm{conjC}$ apply Clozel's arithmeticity theorem
(ibid., Thm.~3.13) to the regular algebraic cusp and
$\sigma\in\operatorname{Aut}(\mathbb C)$. It supplies a cusp with
finite representation the scalar conjugate of $\pi_f$ and coefficient
highest weight at $\sigma\tau$ the old one at $\tau$. Hence its
infinitesimal multiset there is $m_\tau$: the integral weights and
the rational $\rho$ are fixed by $\sigma$. The identification of
the finite factors is additive, $\sigma$-semilinear and commutes
with every finite-adelic operator. Tensoring one finite-factor vector
with a fixed nonzero archimedean vector realizes each finite factor
as a subrepresentation of its cusp space, as required by
$\mathrm{IsFinComp}$. This uses scalar conjugation of the finite
representation itself; no choice of square roots in normalized
Satake parameters enters the conjugation.
% END O72H1

% BEGIN O72H2
\paragraph{O72H2. Counting distinct cusps in group cohomology.}
Fix the compact open $U$, the coefficient $\Xi$, and the finite
set $S$ of pairwise nonisomorphic cohomological cusps with
$U$-fixed vectors. Decompose the finite-adelic right cosets into
rational orbits. On an orbit with representative $h$ the stabilizer
is $\Gamma_h=\mathrm{GL}_n(E)\cap hUh^{-1}$. Functions on that
orbit with the source's simultaneous translation and coefficient
action form the coinduced module from $\Xi|_{\Gamma_h}$
(identify a function by evaluation at $h$ and its translates).
Shapiro's lemma identifies its group cohomology with
$H^\bullet(\Gamma_h,\Xi)$. There are finitely many such orbits
by the adelic class-number theorem. This identifies the source's
sum of ranks, $\mathrm{cohDim}$, with the total cohomology dimension
at level $U$. The symmetric space is contractible, and finite
stabilizers have no higher cohomology over $\mathbb C$, so torsion
in the level does not alter this description.

Use the inclusion of cuspidal cohomology into this cohomology
(Franke, \emph{Harmonic analysis in weighted $L_2$-spaces},
Ann. Sci. ENS 31 (1998)). In the automorphic-forms complex the
cuspidal-support decomposition splits off the cuspidal summand;
Franke's comparison identifies the cohomology of this complex
with that of the locally symmetric space. For each $\pi\in S$,
combine a nonzero class at infinity with a nonzero vector of
$\pi_f^U$. In passing from the full real split center to the
space modulo the positive rational split torus, its Lie algebra
factor must be removed: nonzero cohomology requires its scalar
differential on $\pi\otimes\Xi$ to vanish, by contraction with a
central vector. Its cochain factor is then an exterior algebra with
zero differential, so the relative complex after removing that
factor still has nonzero cohomology. The resulting classes lie in
different direct summands of cuspidal cohomology. Choosing one per
summand proves $|S|\le\mathrm{cohDim}(\top,U,\mathrm{ratC}\,\Xi)$,
the literal $\mathrm{borelC}$.
% END O72H2

% BEGIN O72H3
\paragraph{O72H3. Real cusps, smooth duals and admissibility.}
For a totally real field the real model consists of the finite-level,
$K_\infty$-finite cuspidal automorphic forms with moderate growth and
finite orbit under the archimedean center. Its Capelli polynomials at
each infinite place generate that center by the Harish--Chandra
isomorphism. A cusp subspace is an irreducible module for the finite
adeles, $K_\infty$ and the archimedean Lie algebra. It therefore has
the usual archimedean Harish--Chandra factor and smooth finite factor.
Cofinitely many common Satake multisets give isomorphic local factors
there by the spherical Satake correspondence. Jacquet--Shalika strong
multiplicity one gives an isomorphism of the two global cusps, hence
one intertwining precisely these three actions. This is
$\mathrm{smoR}$.

Twist by the group involution $g\mapsto{}^tg^{-1}$. At unramified
places it inverts the Satake characters, so the result has the same
cofinite parameters as the contragredient. This involution on the
adelic group preserves rational points and interchanges opposite
parabolics, hence preserves the cuspidal space. Strong multiplicity
one identifies the twist with the contragredient. Evaluation against
this isomorphism is a complex bilinear form on two copies of $\pi$;
its radicals vanish since the smooth, $K_\infty$-finite dual
separates vectors (finite types and finite level suffice). At finite
places its invariance uses inverse transpose on the second factor;
at each real orthogonal compact subgroup the involution is the
identity. Its differential is $X\mapsto-{}^tX$, giving the indicated
sum of Lie derivatives equal to zero. This proves $\mathrm{gkR}$.

For $\mathrm{admR}$ choose a compact open $U_0$ inside the specified
open $U$. The finite factor's $U_0$-invariants are finite dimensional
by admissibility. The abstract finite-dimensional input $\xi$ to a
compact intertwiner need not be continuous. It still has finite
composition length. Every irreducible compact type in the image of
any of the tested intertwiners is a composition factor of $\xi$,
since that image is a quotient; the image in the automorphic space
is a finite-dimensional continuous compact representation and is
semisimple. Take the sum of these finitely many continuous isotypic
subspaces in $\pi^{U_0}$. Each is finite dimensional by
Harish--Chandra admissibility at infinity and the finite-adelic bound.
This one space works simultaneously for every such map.

Finally the extra scalar hypothesis on the action of the real identity
matrix fixes the action of the connected positive split center:
integrate the scalar differential along its exponential. Twist by an
adelic determinant norm power to make this action unitary (or trivial).
Cuspidal rapid decay then places the form in the cuspidal Hilbert space
modulo that center. The discrete cuspidal decomposition there is
algebraic on the given vector: it has fixed finite level, finitely
many compact types and finitely many infinitesimal characters by
central finiteness. Automorphic finiteness for this level, list of
types and finite-codimension central annihilator bounds the dimension
of their simultaneous space. Projection to each cuspidal summand
preserves these conditions; the nonzero projections are orthogonal,
so there are only finitely many. They are smooth, compact-finite
vectors in irreducible cuspidal modules. Undoing the twist puts the
original vector in the algebraic sum of the native cusp subspaces,
which gives $\mathrm{ssR}$. At rank zero the adelic group is trivial
and all of the corresponding representation statements reduce to a
one-dimensional space.
% END O72H3

% BEGIN O72H4
\paragraph{O72H4. All three cyclic base-change entries.}
Use Arthur--Clozel, \emph{Simple algebras, base change, and the advanced
theory of the trace formula}, Ch.~3, especially Thms.~4.2 and 4.5,
for prime cyclic base change on $\mathrm{GL}_n$. The cuspidal
lift remains cuspidal unless the original cusp has a nontrivial
self-twist by a character of the cyclic quotient. Global class
field theory realizes that character on the ideles with norm
kernel. At unramified places restriction of Weil parameters sends a
Frobenius element upstairs to the residue-degree power downstairs.
It therefore raises every Satake parameter to that power. At a
complex place over a complex place the infinity parameter is carried
through its chosen embedding.

Conversely, the cofinite equality over places over the same prime
implies that each Galois conjugate of the given cusp has the same
cofinite unramified components. Strong multiplicity one makes it
Galois invariant. Cyclic descent now supplies a cuspidal preimage.
Two cuspidal preimages differ by a character of the cyclic quotient.
This last assertion applies to the source's weak preimages as
well: a noncuspidal base change would have a proper isobaric
decomposition, which cannot have the same cofinite parameters as the
given cusp by isobaric uniqueness. At almost every place the twist
multiplies each parameter by its character on a uniformizer.
Together these give the ascent, descent and fiber clauses of
$\mathrm{cyc}$.

For a quadratic CM extension of a totally real field, use the same
dichotomy in rank two. The associated quadratic character has value
$+1$ at split primes and $-1$ at inert primes; its archimedean signs
have zero differential. A self-twist intertwiner therefore has the
finite, compact and unmodified Lie actions in $\mathrm{IsoTwist}$.
In the cuspidal branch restriction from $W_{\mathbb R}$ to
$W_{\mathbb C}$ preserves the infinitesimal multiset at the
specified embedding; at the finite places the preceding power
identity is precisely the cofinite clause of $\mathrm{IsBC0}$.
This supplies $\mathrm{bc}$. For an imaginary quadratic field over
$\mathbb Q$, apply the descent argument to the given invariant
rank-two cusp. At split primes its Satake multiset is unchanged.
At the complex place the infinitesimal character of the restricted
real parameter agrees at each of the two embeddings with that of
the descent. This yields $\mathrm{DescS}$ under the given regularity
and gap conditions, hence the first alternative of $\mathrm{cyc2}$.
% END O72H4

% BEGIN O72H5
\paragraph{O72H5. The newform normalization.}
Adelize the level-$M$, weight-$k$ holomorphic cusp form with the
specified nebentypus. Its first Fourier coefficient is one, so its
adelization is nonzero. The given coefficient recurrence for every
$p\nmid M$ is the Fourier expansion of the Hecke eigen-equation;
uniqueness of the Fourier expansion gives that equation on the
adelic vector. It has scalar action of the split center in unitary
normalization. By O72H3 choose a nonzero cuspidal component in its
finite decomposition. Projection commutes with all the Hecke
operators and the archimedean center, so this one component works
at every good prime. At infinity its two parameters are
$\pm(k-1)/2$. At $p\nmid M$ it has a spherical vector, and for its
normalized Satake parameters $\alpha,\beta$, the two numbers
$p^{(k-1)/2}\alpha,p^{(k-1)/2}\beta$ have sum $a_p$ and product
$\chi(p)p^{k-1}$. These are the degree-one and central Hecke
operators on that vector.

Twist this cusp by $|\det|^s$ with $s=1-k/2+j$ for the given
integer $j$. The infinity entries shift by $s$ to
$1/2+j$ and $3/2-k+j$. The finite Satake parameters are multiplied
by $p^{-s}$. In rank two the native $\mathrm{algSat}$ supplies
$p^{1/2}$; multiplying also by $p^j$ gives precisely
$p^{j+1/2-s}=p^{(k-1)/2}$. The product of the two linear factors
in the contract is consequently $X^2-a_pX+\chi(p)p^{k-1}$ for
every $p\nmid M$. This gives $\mathrm{newform}$ with the required
twist, including negative $j$.
% END O72H5

% BEGIN O72H6
\paragraph{O72H6. Rank-two attachment over the original totally real field.}
Let $n=2$, fix the prime $\ell$, the coefficient isomorphism $\iota$,
and the given regular algebraic polarized real cusp $\pi$. Its
nondegenerate multiplier pairing identifies its dual with the
corresponding character twist: for a finite-level compact-finite
vector the functional in the second argument belongs to the
admissible dual, and the nonzero intertwiner between irreducibles
is an isomorphism. The supplied character has constant sign on
the real places. Its infinity exponents are integral by the
regular algebraic infinitesimal characters: the pair of entries at
each real place is the negative pair translated by the multiplier
exponent, which is their sum in rank two. This is the polarized
regular algebraic hypothesis of Barnet-Lamb--Gee--Geraghty--Taylor
(BLGGT), \emph{Potential automorphy and change of weight},
Thm.~2.1.1. The theorem gives a continuous semisimple $r$ on the
specified Galois group (transport through an identification of
algebraic closures), de Rham at $\ell$, with integral jumps and
a common purity weight $w_0$. Choose a basis in its finite field
of definition before extension to $\overline{\mathbb Q}_\ell$.
The Harish--Chandra formula in this convention reads
$\mathrm{HT}_\tau=\{1/2-z:z\in m_{\iota\tau}\}$: the shift
by $\rho$ in the dual coefficient and the shift by $(n-1)/2$
in the correspondence give the usual jumps
$\lambda_1+1,\lambda_2$ for the coefficient of highest weight
$(\lambda_1,\lambda_2)$. The symmetry in that theorem pairs
$h$ with $w_0-h$ at the same embedding over a totally real field.

At $v\nmid\ell$ the theorem identifies the Frobenius-semisimplified
Weil--Deligne representation with the parameter of
$\pi_v|\det|^{-1/2}$ and gives its purity. A spherical factor of
a global cusp is generic; its unramified character support is
unlinked (an adjacent pair would make its spherical constituent
nongeneric). Hence it is plain in the sense of the source. Its
parameter has trivial inertia and zero monodromy, so $r$ is
unramified there. At the inverse of arithmetic Frobenius the
eigenvalues under $\iota$ are $q^{1/2}$ times the unitary-normalized
Satake entries, exactly $\mathrm{FrobEig}$. At any plain toral
place the factor is induced from unlinked rank-one characters.
Its normalized Langlands parameter is their direct sum after the
indicated twist, with zero monodromy. At such a place away from
$\ell$, purity implies every degree-one Weil element has
eigenvalues all of the same absolute value. Indeed inertia has
finite image and a suitable power of the element is a power of
Frobenius acting trivially on that inertia by conjugation; purity
with $N=0$ gives the common absolute value, unchanged by taking
the corresponding roots. Since $q>1$ and the eigenvalues are
nonzero this gives the strict inequality in $\mathrm{RatioW}$.

For $v\mid\ell$ use BLGGT, \emph{Local-global compatibility for
$l=p$, II}, Thm.~A. Its semisimplified Weil comparison applies to
totally real fields with ramified local factors as well. In the
plain toral case the two Weil characters have no relative
norm twist, since they came from unlinked rank-one support.
The relation $r_W(g)Nr_W(g)^{-1}=|g|N$ forces $N=0$: otherwise
there would be a common constituent of the semisimplified
Weil representation and its norm twist linked in this way.
Thus we recover equality of Frobenius-semisimplified
Weil--Deligne parameters on the full plain range, including
ramified characters. At a spherical place, crystallinity follows
directly from Thm.~2.1.1(4).

Finally the local comparison in this record tests every
$\mathrm{IsRecT}$, not just one chosen correspondence. On a toral
unlinked induction its value is fixed by its rank-one and
Langlands-sum laws. Its rank-one map uses the usual local Artin
normalization: the global compatibility field tests a normalized
global Artin map, and its restriction to one local multiplicative
group is injective. For the valuation use the norm to $\mathbb Q$
and cyclotomic characters away from the residue prime to detect
any nonzero valuation. For a unit in the kernel, approximate its
finite idèle by principal idèles in the kernel description. These
are eventually units everywhere. For each $m$, V72L7 gives a modulus
prime to the residue prime forcing such approximating units, which
are close to one away from the chosen place, to be global $m$-th
powers. The local component belongs to every $m$-th power subgroup
of the compact local units, whose intersection is trivial by their
finite quotients. This identifies the
two choices, including the inverse-Frobenius convention.
For the specified local bridges above $\ell$, choose a common finite
extension containing the field of the matrices, the given bridge
coefficient field and the images of the embeddings under test.
Base extension of de Rham invariants and their filtration gives
the specified multiplicities; functoriality for the identification
of the completed local fields gives the same potentially
semistable parameter and crystalline comparison in that model.
F71A3 identifies the canonical period rings and filtrations.
These checks give the local comparison, jumps, Frobenius, ratio
and crystalline clauses on the same $r$. Together they supply
$\mathrm{galR}$.
% END O72H6
